\section{Volume and measure consequences}\label{sec:consequences}

With Theorem~\ref{thm:main} proved, we derive the additive volume
inequality stated in the introduction and then apply Saroglou's
transfer theorem to even log-concave measures.

\subsection{The symmetric \texorpdfstring{$L_p$}{Lp} volume inequality}
\label{subsec:lp-proof}

\begin{proof}[Proof of Corollary~\ref{cor:lp}]
At $\lambda=0,1$, the assertion follows from
$\mathcal W[h_K]=K$ and $\mathcal W[h_L]=L$. Suppose $0<\lambda<1$.
Set
\[
 a=|K|^{1/n}>0,\qquad b=|L|^{1/n}>0,\qquad
 K_0=a^{-1}K,\qquad L_0=b^{-1}L,
\]
so that $|K_0|=|L_0|=1$. Define
\[
 c=\bigl((1-\lambda)a^p+\lambda b^p\bigr)^{1/p}>0,
 \qquad \theta=\frac{\lambda b^p}{c^p}\in(0,1).
\]
Then $1-\theta=(1-\lambda)a^p/c^p$. For every $u\in S^{n-1}$,
the support identities $h_K=a h_{K_0}$ and $h_L=b h_{L_0}$, followed
by the weighted arithmetic--geometric mean inequality, give
\begin{align*}
 \bigl((1-\lambda)h_K(u)^p+\lambda h_L(u)^p\bigr)^{1/p}
 &=c\bigl((1-\theta)h_{K_0}(u)^p+
              \theta h_{L_0}(u)^p\bigr)^{1/p}\\
 &\ge c h_{K_0}(u)^{1-\theta}h_{L_0}(u)^\theta.
\end{align*}
All support values here are positive. The Wulff definition is monotone
in its support bound and satisfies $\mathcal W[cf]=c\mathcal W[f]$ for
$c>0$. Hence
\[
 c\mathcal W[h_{K_0}^{1-\theta}h_{L_0}^\theta]
 \subset
 \mathcal W\!\left[
  \bigl((1-\lambda)h_K^p+\lambda h_L^p\bigr)^{1/p}
 \right].
\]
By Theorem~\ref{thm:main}, the unscaled body on the left has volume at
least $1$, because $K_0,L_0$ are origin-symmetric and have unit volume.
The body on the right therefore has volume at least $c^n$. Raising this
bound to the positive power $p/n$ gives
\eqref{eq:lp}, since $c^p=(1-\lambda)|K|^{p/n}+\lambda|L|^{p/n}$.
\end{proof}

\subsection{Even log-concave measures and the \texorpdfstring{$(B)$}{(B)}-conjecture}
\label{sec:measure}

We now pass from volume to integration against even log-concave measures.
A nonnegative Radon measure $\mu$ on $\R^n$ is \emph{log-concave} if
\[
 \mu((1-\theta)A+\theta B)
 \ge \mu(A)^{1-\theta}\mu(B)^\theta
 \qquad(0<\theta<1)
\]
for all compact sets $A,B\subset\R^n$; it is \emph{even} if
$\mu(A)=\mu(-A)$ for every Borel set $A$. Radon measures are finite on
compact sets, but their total mass need not be finite. We use the analogous
multiplicative definition of log-concavity for nonnegative functions, so
the identically zero function is included. Endpoint products in an
interpolation inequality are understood as the corresponding endpoint
values.

Saroglou's transfer theorem states that the logarithmic Brunn--Minkowski
inequality for Lebesgue measure in a dimension $d$ implies the same
inequality in dimension $d$ for every even log-concave density
\cite[Theorem~3.1]{Saroglou2016}. The density may vanish and need not be
normalized. This theorem is stated for densities with respect to
$d$-dimensional Lebesgue measure. The reduction below makes explicit how
it also gives the measure statement when the measure has lower-dimensional
support.

\begin{corollary}[Measure inequality and scalar $(B)$-conjecture]
\label{cor:measure}
Let $n\ge1$, let $\mu$ be an even log-concave Radon measure on $\R^n$, and
let $K,L\subset\R^n$ be origin-symmetric convex bodies. Then, for every
$0\le\lambda\le1$,
\begin{equation}\label{eq:measure-logbm}
 \mu\bigl(\mathcal W[h_K^{1-\lambda}h_L^\lambda]\bigr)
 \ge \mu(K)^{1-\lambda}\mu(L)^\lambda.
\end{equation}
Consequently, for every origin-symmetric convex body $K\subset\R^n$, the
function
\begin{equation}\label{eq:b-function}
 \R\ni t\longmapsto \mu(e^tK)
\end{equation}
is log-concave.
\end{corollary}

\begin{proof}
Fix $0<\lambda<1$ and write
$D=\mathcal W[h_K^{1-\lambda}h_L^\lambda]$.
The assertion is immediate when $\mu$ is the zero measure. Otherwise,
choose a closed ball $C$ centered at the origin large enough that
$K\cup L\cup D\subset C$ and $0<\mu(C)<\infty$. This is possible because
the three bodies are compact and $\mu$ is a nonzero Radon measure. The
probability measure
\[
 \nu(A)=\frac{\mu(A\cap C)}{\mu(C)}
\]
is even and log-concave: for compact $A,B$, convexity of $C$ gives
\[
 (1-\theta)(A\cap C)+\theta(B\cap C)
 \subset ((1-\theta)A+\theta B)\cap C,
\]
and the defining inequality for $\mu$ applies. Since $K,L,D\subset C$,
the common factor $\mu(C)$ cancels from \eqref{eq:measure-logbm}. It
therefore suffices to prove that inequality for $\nu$.

Let $E$ be the affine hull of $\operatorname{supp}\nu$. Evenness gives
$E=-E$, so $E$ is a linear subspace. If $\dim E=0$, then
$\nu=\delta_0$; all three bodies contain $0$, and equality holds.
Suppose $d=\dim E\ge1$ and identify $E$ isometrically with $\R^d$.
By inner regularity, the compact-set definition
of log-concavity for $\nu$ is equivalent to the Borel-set formulation
using inner measure. Borell's characterization
\cite[Theorem~1.1, with $s=0$]{Borell1974}, applied to $\nu$ on $E$,
represents it as an affine image of a log-concave density in a dimension
$k\le d$. Since its support spans $E$, the affine map has rank $d$.
Thus $k=d$, the map is invertible, and $\nu$ has a log-concave density
$f$ with respect to Lebesgue measure on $E$. Evenness of $\nu$ gives
$f(x)=f(-x)$ almost everywhere; replacing $f(x)$ by
$\sqrt{f(x)f(-x)}$ gives an even log-concave representative of the same
density.

Set $K_E=K\cap E$ and $L_E=L\cap E$. These are origin-symmetric convex
bodies in $E$, since $K$ and $L$ contain neighborhoods of $0$. Define
their relative support functions for every $v\in E$ by
\[
 h^E_{K_E}(v)=\max_{y\in K_E}\langle y,v\rangle,
 \qquad
 h^E_{L_E}(v)=\max_{y\in L_E}\langle y,v\rangle.
\]
The relative Wulff body in $E$ is
\[
 D_E=\bigcap_{v\in E}
 \left\{x\in E:\langle x,v\rangle
 \le h^E_{K_E}(v)^{1-\lambda}h^E_{L_E}(v)^\lambda\right\}.
\]
Using all $v\in E$ rather than only unit vectors gives the same body by
homogeneity. Let $P_E$ be the orthogonal projection onto $E$. For
$x\in D_E$ and $u\in S^{n-1}$,
\begin{align*}
 \langle x,u\rangle
 &=\langle x,P_Eu\rangle\\
 &\le h^E_{K_E}(P_Eu)^{1-\lambda}
       h^E_{L_E}(P_Eu)^\lambda\\
 &\le h_K(u)^{1-\lambda}h_L(u)^\lambda.
\end{align*}
Indeed, $h^E_{K_E}(P_Eu)=\max_{y\in K\cap E}\langle y,u\rangle\le h_K(u)$,
and likewise for $L$; if $P_Eu=0$, the middle product is $0$.
Thus
\begin{equation}\label{eq:relative-wulff-inclusion}
 D_E\subset D\cap E.
\end{equation}
Theorem~\ref{thm:main} in dimension $d$, followed by Saroglou's transfer
theorem \cite[Theorem~3.1]{Saroglou2016} in that same dimension, gives
\[
 \nu(D_E)\ge \nu(K_E)^{1-\lambda}\nu(L_E)^\lambda.
\]
Since $\nu$ is supported on $E$, inclusion
\eqref{eq:relative-wulff-inclusion} yields
\[
 \nu(D)=\nu(D\cap E)\ge\nu(D_E)
 \ge\nu(K)^{1-\lambda}\nu(L)^\lambda,
\]
as required. At $\lambda=0,1$, the Wulff bodies are $K,L$, respectively,
and the endpoint convention gives equality.

Finally, for $s,t\in\R$, the support identity $h_{e^sK}=e^s h_K$ gives
\[
 \mathcal W\!\left[h_{e^sK}^{1-\lambda}h_{e^tK}^\lambda\right]
 =\mathcal W\!\left[e^{(1-\lambda)s+\lambda t}h_K\right]
 =e^{(1-\lambda)s+\lambda t}K.
\]
Applying \eqref{eq:measure-logbm} to $e^sK$ and $e^tK$ proves
\[
 \mu\bigl(e^{(1-\lambda)s+\lambda t}K\bigr)
 \ge \mu(e^sK)^{1-\lambda}\mu(e^tK)^\lambda.
\]
This is exactly the log-concavity in \eqref{eq:b-function}, and is the
scalar-dilation argument of Saroglou \cite[Corollary~3.2]{Saroglou2016}.
\end{proof}
